\documentclass[11pt]{article}
\usepackage[a4paper,margin=28mm]{geometry}
\usepackage{amsmath,amssymb,amsthm,mathtools,hyperref}
\newcommand{\ddc}{dd^c}
\newcommand{\MA}{\operatorname{MA}}
\newcommand{\PSH}{\operatorname{PSH}}
\newcommand{\vol}{\operatorname{vol}}
\newcommand{\rank}{\operatorname{rank}}
\newcommand{\1}{\mathbf 1}
\newtheorem{theorem}{Theorem}
\newtheorem{lemma}[theorem]{Lemma}
\newtheorem{proposition}[theorem]{Proposition}
\newtheorem{example}[theorem]{Example}
\theoremstyle{remark}\newtheorem{remark}[theorem]{Remark}
\title{Smaller-background obstacle envelopes and smooth extension\newline under dual attainment}
\author{}
\date{}
\begin{document}
\maketitle
\begin{abstract}
We give two sufficient conditions for a genuine smaller-background obstacle
envelope: constant bounded Monge--Amp\`ere mass on the specified Hermitian
background, and a quantitative sublinear lower bound on singularities.
Both conditions provide actual canonical competitors at every smaller scale
and force zero Lelong numbers. Under attained duality, a globally smooth
semipositive root of complex rank at least two then gives smooth extension of
the threshold potential across its analytic degeneracy set. A Hopf surface
and a quadric double cover verify actual attainment. An unbounded projective
example has a full-mass density in no $L^p$, $p>1$, but cannot have the
specified smooth maximizer. The general moment and entropy implications
remain unresolved.
\end{abstract}

\section{Conventions and the precise scope}
Let $X$ be a connected compact complex manifold of dimension $n\geq2$,
$h$ a smooth positive Hermitian form, and $dV$ a smooth positive probability.
We use $\ddc=\frac{i}{2\pi}\partial\bar\partial$ and ordinary exterior
powers. Potentials are genuine canonical quasi-plurisubharmonic functions;
none is changed on a null set. Write
\[
 \MA_h^{\rm ext}(u)=\lim_{j\to\infty}
 \1_{\{u>-j\}}(h+\ddc\max(u,-j))^n
\]
for normalized $h$-psh $u\leq0$, using the non-pluripolar truncation
convention. Initially this measure may have infinite total mass.

For nonzero $f\geq0$ in $L^1(dV)$, $b_h(f)$ is the maximal scalar $b$
for which a genuine potential has $\MA_h^{\rm ext}(u)\geq bf\,dV$.
We use the supplied threshold existence theorem, its greatest normalized
member $\Phi_f$, and its characterization by decreasing bounded-equation
normalizing scalars for $f_k=\min(f,k)$.
A smooth positive Gauduchon test is an $(n-1,n-1)$ form $T$ with
$\ddc T=0$ and $\int_Xh\wedge T=1$.
For $T=\gamma^{n-1}$ and $p=f^{1/n}$ put
\[
 k_T=\left(\frac{\gamma^n}{dV}\right)^{(n-1)/n},\qquad
 D_p(T)=\int_Xp k_T\,dV.
\]
The supplied dual formula is $\sup_T D_p(T)=b_h(f)^{-1/n}$.
When this supremum is attained by a smooth semipositive $T$ positive off a
proper analytic $Z$, we also use the supplied exact finite threshold equation,
smoothness of $\Phi_f$ off $Z$, and uniqueness up to constants.
These are background assumptions, not conclusions proved here.

For $0<\delta<1$ the obstacle envelope is
\[
 P_{\delta h}(u)=\left(\sup\{v\in\PSH(X,\delta h):v\leq u\}\right)^*.
\]
It is genuine only when its competitor family is nonempty and its canonical
envelope is not identically $-\infty$. This is an obstacle envelope, with
pointwise canonical inequalities, rather than a zero-obstacle envelope of
a singularity class. The known capacity-convergence consequence of a genuine
smaller-background envelope applies to our criteria. We make no assertion
about smoothness of such an envelope at its free boundary.

The bounded Hermitian product, contact comparison and envelope-support
premises used below are those of \cite[\S\S1--2]{GL}.
They are stated for positive, possibly nonclosed backgrounds.
The constant-mass metric condition is classical
\cite[Proposition 3.3(4) and its discussion]{GL}; no priority claim is made.

\section{A geometric condition on the specified background}
\begin{theorem}\label{th:mass}
Suppose
\begin{equation}\label{eq:G}
 \ddc(h^k)=0\qquad(1\leq k<n).
\end{equation}
Set $V=\int_Xh^n$ and $M=\int_Xf\,dV$. Then $b_h(f)=V/M$ and every
threshold subsolution solves $\MA_h^{\rm ext}(u)=b_h(f)f\,dV$, of mass $V$.
For every genuine $h$-psh potential of full mass $V$,
$P_{\delta h}(u)$ is genuine for every $0<\delta<1$.
Every such potential has zero Lelong numbers.
Dual attainment and boundedness of the potential are unnecessary.
\end{theorem}
\begin{proof}
For a bounded $h$-psh $w$, the bounded-product expansion and weighted Stokes give
\[
 \int_X(\ddc w)^\ell\wedge h^{n-\ell}
 =\int_Xw(\ddc w)^{\ell-1}\wedge\ddc(h^{n-\ell})=0
 \quad(1\leq\ell\leq n).
\]
For $\ell=n$ the last factor is $\ddc1=0$.
Thus $\int_X\MA_h(w)=V$. For an unbounded normalized $w$ the truncation
definition gives $\MA_h^{\rm ext}(w)(X)\leq V$.
The supplied bounded equations have normalizers
$c_k=V/\int_Xf_k\,dV$, so their decreasing limit is $V/M$.
A threshold subsolution has mass at least $(V/M)M=V$ and at most $V$;
the nonnegative difference of its measure and $(V/M)f\,dV$ therefore vanishes.

Now let $u\leq0$ be any full-mass potential, with
$\mu=\MA_h^{\rm ext}(u)$, and fix $a>1$.
Put $u_j=\max(u,-j)$ and $w_j=P_h(au_j)$.
The constant $-aj$ is a competitor, so $w_j$ is bounded between $-aj$ and $0$.
The obstacle is bounded and quasi-continuous. Theorem 2.3 of \cite{GL}
places its measure on $C_j=\{w_j=au_j\}$.
The source defines envelopes by quasi-everywhere inequalities. Here both
the obstacle and any competitor become psh after adding a common local
smooth potential with sufficiently large Hessian. A quasi-everywhere
inequality is then an almost-everywhere inequality between psh functions,
and hence a canonical pointwise inequality. The two definitions coincide
for these particular obstacles. Upper semicontinuity of $au_j$ also ensures
$w_j\leq au_j$ everywhere.

Both $w_j$ and $au_j$ are $ah$-psh. Positivity of bounded products and
contact comparison \cite[Lemma 1.2]{GL} give
\begin{equation}\label{eq:contact}
 \MA_h(w_j)
 \leq\1_{C_j}\MA_{ah}(w_j)
 \leq\1_{C_j}\MA_{ah}(au_j)
 =a^n\1_{C_j}\MA_h(u_j).
\end{equation}
For integers $j\geq N$, locality and the exact bounded mass imply
\begin{equation}\label{eq:tail}
 \int_{\{u_j\leq-N\}}\MA_h(u_j)
 =V-\mu(\{u>-N\})=\mu(\{u\leq-N\}).
\end{equation}
The last quantity tends to zero because $\mu$ is finite and gives no mass
to the pole set. Choose $N$ with $a^n\mu(\{u\leq-N\})<V$.
If $\sup_Xw_j<-aN$, its contact set lies in $\{u_j<-N\}$, and
\eqref{eq:contact}--\eqref{eq:tail} contradict $\int_X\MA_h(w_j)=V$.
Consequently $\sup_Xw_j\geq-aN$; for $j<N$ this follows from $w_j\geq-aj$.
The decreasing limit $w$ is canonical $h$-psh and genuine: a decreasing
sequence of quasi-psh functions with uniformly bounded Hessian below cannot
collapse to $-\infty$ while its suprema stay bounded below.
Pointwise $w\leq au$, including at poles.
For $\delta=1/a$, the function $v_\delta=\delta w$ is a genuine
$\delta h$-psh competitor below $u$. The envelope is bounded below by this
competitor and above by the upper semicontinuous $u$, so is genuine.
The zero-Lelong assertion follows from the next lemma.
\end{proof}

\begin{lemma}\label{lem:lelong}
If a genuine $h$-psh $u$ has a genuine $\delta h$-psh competitor below it
for every $0<\delta<1$, then $u$ has zero Lelong numbers everywhere.
\end{lemma}
\begin{proof}
Fix a smooth strictly positive normalized Gauduchon test $G$.
Such a test exists on every compact complex manifold
\cite[Theorem 2, pp.~132--133]{Gaud}; it is independent of the
possibly degenerate maximizer.
For any such competitor $v_\delta$,
\[
 \int_X(\delta h+\ddc v_\delta)\wedge G=\delta.
\]
On a fixed coordinate ball choose a closed Euclidean form $\omega_0$ and
$K$ with $h\leq K\omega_0$. The closed positive current
$B_\delta=\ddc v_\delta+\delta K\omega_0$ has local trace mass at most
$C\delta$: the strictly positive $G$ controls the trace of
$\delta h+\ddc v_\delta$, and the remaining correction is smooth of size
$O(\delta)$. The local Lelong mass estimate on a smaller ball gives
$\nu(v_\delta,x)\leq C\delta$. A finite atlas makes $C$ uniform.
Since $v_\delta\leq u$, Lelong comparison gives
$0\leq\nu(u,x)\leq\nu(v_\delta,x)\leq C\delta$.
Let $\delta$ decrease to zero.
\end{proof}

\begin{example}\label{ex:hopf}
On $X=(\mathbb C^2\setminus\{0\})/\langle z\mapsto2z\rangle$,
put $s=|z|^2$ and $h=s^{-1}\ddc s$. It descends and is positive.
For $F(s)=s^{-1}$, $sF''+2F'=0$, whence
$\ddc h=\ddc(s^{-1})\wedge\ddc s=0$, whereas $dh\ne0$.
The quotient is diffeomorphic to $S^3\times S^1$, so is non-K\"ahler.
Choose nonconstant smooth real $q$ and small nonzero $\epsilon$ so that
$A=h+\ddc(\epsilon q)>0$. Set $u=\epsilon q-\sup_X\epsilon q$,
$dV=h^2/V$, $f=A^2/dV$, and $T=A/V$.
Theorem~\ref{th:mass} gives $b=1$. Further,
$\ddc T=0$, $\int_Xh\wedge T=1$, and $D_{\sqrt f}(T)=1$.
For every normalized positive Gauduchon test $S$, matrix AM--GM and Stokes
give $D_{\sqrt f}(S)\leq\int_XA\wedge S=1$.
Thus $T$ really attains the supremum and supplied uniqueness identifies
$\Phi_f=u$. This is a nonconstant smooth attained instance, not an
unbounded example. On any surface, \eqref{eq:G} means $\ddc h=0$ on the
specified background itself.
\end{example}

\section{Across-set extension from a smooth root of rank at least two}
\begin{theorem}\label{th:extension}
Retain the attained setting from Section 1. Suppose $T=\gamma^{n-1}$ for
a globally smooth semipositive $(1,1)$ form $\gamma$ with
$\rank_{\mathbb C}\gamma_x\geq2$ at every point, and suppose $\Phi_f$
has zero Lelong numbers. Then $\Phi_f$ is smooth on all $X$.
In particular \eqref{eq:G} supplies the zero-Lelong hypothesis.
\end{theorem}
\begin{proof}
Set $U=X\setminus Z$ and $A=h+\ddc\Phi_f$.
Distributional Stokes gives $\int_XA\wedge T=1$.
On $U$ the forms are smooth, $\gamma>0$, and AM--GM gives
\[
 A\wedge\gamma^{n-1}\geq(A^n)^{1/n}(\gamma^n)^{(n-1)/n}.
\]
The integral of the right side is $b^{1/n}D_p(T)=1$ by the exact equation
and attainment; $Z$ has smooth-volume measure zero.
The continuous nonnegative gap therefore vanishes on $U$.
Equality gives $A=r\gamma$ for smooth $r\geq0$ there.
If $\det A=0$, equality forces $\operatorname{tr}_\gamma A=0$ and $A=0$,
so $r=\operatorname{tr}_\gamma A/n$ is still smooth.

Since $\ddc\Phi_f$ is closed,
\begin{equation}\label{eq:scalar}
 dr\wedge\gamma+r\,d\gamma=dh\qquad\text{on }U.
\end{equation}
For a semipositive real $(1,1)$ form of complex rank at least two the map
$W_\gamma:\Lambda^1_{\mathbb R}\to\Lambda^3_{\mathbb R}$,
$\xi\mapsto\xi\wedge\gamma$, is injective.
Indeed, diagonalize into disjoint positive coordinate two-planes.
For each coordinate of $\xi$, choose a positive plane not containing it.
The corresponding triple coefficient is that coordinate times a positive
eigenvalue and cannot receive a contribution from another disjoint pair.
All coordinates therefore vanish if the wedge vanishes.
With an auxiliary metric its left inverse
$L_\gamma=(W_\gamma^*W_\gamma)^{-1}W_\gamma^*$ is smooth and uniformly
bounded on compact $X$. Equation~\eqref{eq:scalar} implies
\begin{equation}\label{eq:ode}
 dr=\alpha-r\beta,\qquad
 \alpha=L_\gamma(dh),\quad\beta=L_\gamma(d\gamma).
\end{equation}
In particular $|d\log(1+r)|\leq C$ on $U$.

We spell out why the analytic set causes no path obstruction.
Locally it has a countable smooth stratification of real dimensions at most
$2n-2$. For an outside point, the endpoints of segments meeting a stratum
lie in cones of dimension at most $2n-1$, hence a null set.
Two outside points in a coordinate ball can be joined through a generic
third point by two segments avoiding $Z$. Choosing a ball of radius
comparable to their distance makes the length at most a fixed multiple of
that distance, after shrinking the chart. A finite connected coordinate
cover and outside points in its overlaps give a global uniform bound on
connecting-path lengths in $U$.
The logarithmic gradient bound consequently bounds $r$ on $U$.
Equation~\eqref{eq:ode} then bounds $dr$, and the short paths give a unique
locally Lipschitz extension of $r$ across $Z$.
As $Z$ is Lebesgue-null, \eqref{eq:ode} holds for its weak derivative on all
$X$. The right side is continuous, so $r$ is $C^1$; iteration makes it smooth.

Put $S=r\gamma$. It is smooth semipositive and $dS=dh$ everywhere by density.
The current $R=A-S$ is closed and supported on $Z$.
It is positive: outside $Z$, $A=S$, while $S$ has no mass on $Z$;
thus $R$ is exactly the positive matrix-measure part of $A$ on $Z$.
The analytic support theorems \cite[III.2.10--2.14]{Dem} give
$R=\sum_i c_i[D_i]$, $c_i\geq0$, on the divisorial components;
components of codimension at least two carry no such current.
Locally $S-h$ is smooth and closed and has a smooth $\ddc$ potential.
Thus the Lelong numbers of $R$ equal those of $\Phi_f$.
They are zero, so every $c_i=0$ and $\ddc\Phi_f=S-h$ globally.
Taking a local Euclidean trace gives an elliptic distributional equation
with smooth right side. Its solution is smooth, and its canonical
quasi-psh representative agrees with that smooth solution everywhere.
\end{proof}

\begin{example}\label{ex:quadric}
For $n\geq3$ let
\[
 X=\{[w:z_0:\cdots:z_n]\in\mathbb P^{n+1}:w^2=\textstyle\sum_i z_i^2\},
 \qquad F([w:z])=[z].
\]
The defining quadratic has no projective singular point. This smooth
connected compact quadric double-covers $\mathbb P^n$ and has smooth
ramification divisor $Z=\{w=0\}$.
Take restricted normalized Fubini--Study $h$ and $A=F^*\omega_{\rm FS}$.
The global smooth function
\[
 u=\log\frac{\sum_i|z_i|^2}{|w|^2+\sum_i|z_i|^2}
\]
satisfies $A=h+\ddc u$, $-\log2\leq u\leq0$, and $\sup_Xu=0$ on $Z$.
The differential of $F$ has rank $n$ off $Z$ and $n-1$ along $Z$.
With $V=\int_Xh^n=\int_XA^n=2$, put $dV=h^n/V$, $f=A^n/dV$,
$T=A^{n-1}/V$ and $\gamma=A/V^{1/(n-1)}$.
Stokes gives $\int_Xh\wedge T=1$ and $\ddc T=0$.
Also $D_p(T)=1$ and Theorem~\ref{th:mass} gives $b=1$.
For all normalized positive tests $S$, $D_p(S)\leq\int_XA\wedge S=1$;
the semipositive $A$ is allowed in this matrix inequality.
Thus $T$ is an actual smooth maximizer, positive off $Z$, and $\Phi_f=u$.
Its smooth root has rank at least two everywhere and truly degenerates on
$Z$. In dimension two the rank hypothesis of Theorem~\ref{th:extension}
instead forces positivity everywhere.
\end{example}

\section{A sublinear singularity criterion}
\begin{theorem}\label{th:sublinear}
For arbitrary $h$, suppose a genuine global quasi-psh $\psi\leq0$ satisfies
$\ddc\psi\geq-Ch$, $C>0$, and
\begin{equation}\label{eq:sublinear}
 \Phi_f\geq-A(1-\psi)^\alpha-B,
 \qquad A>0,\ B\in\mathbb R,\ 0<\alpha<1,
\end{equation}
canonically pointwise, with right side $-\infty$ at poles of $\psi$.
Then every $P_{\delta h}(\Phi_f)$ is genuine and $\Phi_f$ has zero Lelong
numbers. Together with the smooth-root and rank hypotheses of
Theorem~\ref{th:extension}, this gives smooth extension across $Z$.
\end{theorem}
\begin{proof}
Choose $\alpha<\beta<1$ and $\epsilon>0$ with $\epsilon\beta C\leq\delta$.
The finite constant
\[
 K=B+\sup_{t\geq1}(At^\alpha-\epsilon t^\beta)
\]
makes $v_\delta=-\epsilon(1-\psi)^\beta-K\leq\Phi_f$ pointwise.
The composition is increasing and convex, with
$0<\chi'(t)=\epsilon\beta(1-t)^{\beta-1}\leq\epsilon\beta$ and
$\chi''(t)=\epsilon\beta(1-\beta)(1-t)^{\beta-2}\geq0$.
Apply the quasi-psh chain inequality to $\max(\psi,-j)$ and pass to the
decreasing limit: $\ddc v_\delta\geq-\epsilon\beta Ch\geq-\delta h$.
It is genuine because $\psi$ is finite almost everywhere.
At a pole its value is canonically $-\infty$, and the comparison remains
valid. Thus it is an actual competitor for every $\delta$.
Lemma~\ref{lem:lelong} and Theorem~\ref{th:extension} finish the claims.
\end{proof}
For example an effective divisor with a smooth bundle metric supplies
$\psi=\log|s_D|_H^2$, shifted to have supremum at most zero;
Poincar\'e--Lelong gives $\ddc\psi=[D]-\Theta_H\geq-Ch$.
The quantitative bound \eqref{eq:sublinear} is a separate hypothesis.
It is neither boundedness nor an assertion of the desired envelope.

\section{An unbounded example and its missing attainment}
\begin{proposition}\label{prop:excluded}
On $\mathbb P^n$, $n\geq2$, take $h=\omega_{\rm FS}$ with
$\int h^n=1$ and $dV=h^n$. For $0<\alpha<1$ set
\[
 \psi=\log\frac{\sum_{i=1}^n|z_i|^2}{\sum_{i=0}^n|z_i|^2},\qquad
 u=1-(1-\psi)^\alpha.
\]
This genuine normalized potential has sole pole $o=[1:0:\cdots:0]$,
full non-pluripolar mass one, and genuine smaller-background envelopes at
every scale. Its density $f=\MA_h^{\rm ext}(u)/dV$ is in $L^1$ and in
no $L^p$ with $p>1$. Moreover $b_h(f)=1$ and $\Phi_f=u$.
It has no smooth semipositive Gauduchon maximizer positive off a proper
analytic set, so is not a counterexample in the original attained setting.
\end{proposition}
\begin{proof}
Locally $h+\ddc\psi=\ddc\log\sum_{i=1}^n|z_i|^2\geq0$.
The increasing convex composition has derivative at most $\alpha$,
so $u$ is genuine and $h+\ddc u\geq(1-\alpha)h$ off $o$.
Its supremum is zero and Theorem~\ref{th:sublinear}'s competitor construction
applies to $u$, with $A=1$, $B=-1$.

Full mass requires an additional calculation.
In the chart $z_0=1$, $s=\sum_{i=1}^n|z_i|^2$, a local potential is
\[
 F(s)=\log(1+s)+1-\left(1-\log\frac{s}{1+s}\right)^\alpha.
\]
Its radial flux is
\[
 q(s)=sF'(s)=\frac{s+\alpha(1-\log(s/(1+s)))^{\alpha-1}}{1+s},
 \qquad q(0)=0.
\]
Radial Stokes with our normalization gives mass
$q(b)^n-q(a)^n$ on $\{a<|z|^2<b\}$.
For $u_j=\max(u,-j)$, the total mass is one by Theorem~\ref{th:mass}.
On a fixed small ball $\{|z|^2<b\}$, its boundary agrees smoothly with $u$
for large $j$, while its central potential is smooth with zero central flux.
Bounded-product Stokes, or a decreasing local smoothing of the truncation,
therefore gives total ball mass $q(b)^n$.
The non-pluripolar measure of $u$ has the same ball mass by letting $a\downarrow0$.
Outside that ball $u_j$ eventually equals $u$.
Consequently its full non-pluripolar mass is exactly one, without inferring
full mass from zero Lelong numbers alone.

Write $L=1-\log s$. Near $o$ the tangential eigenvalues are asymptotic to
$\alpha s^{-1}L^{\alpha-1}$ and the radial eigenvalue to
$\alpha(1-\alpha)s^{-1}L^{\alpha-2}$.
Thus $f(s)\asymp s^{-n}L^{n(\alpha-1)-1}$.
The radial smooth volume factor is comparable to $s^{n-1}ds$.
The $L^1$ tail is $\int^\infty L^{-1-n(1-\alpha)}dL<\infty$;
every $L^p$ tail, $p>1$, has an exponentially increasing factor
$e^{n(p-1)L}$ and diverges.
The mass bound and this solution give $b=1$.
Full-mass uniqueness in the K\"ahler class
\cite[Theorem 3.1 and \S3.3]{BEGZ} identifies $\Phi_f=u$.

Suppose the specified smooth maximizer $T$ existed.
Integrated AM--GM equality, applied to the smooth positive
$A=h+\ddc u$ off $o$, gives $T=tA^{n-1}$ on the dense set where $T>0$.
By continuity the relation extends on $X\setminus\{o\}$, with smooth
$t\geq0$ recoverable from $(T\wedge A)/A^n$.
Since $A$ is closed, $\ddc T=0$ implies
$\ddc t\wedge A^{n-1}=0$, a strictly elliptic scalar equation there.
All eigenvalues of $A$ tend to infinity at $o$, as do the eigenvalues of
its complementary-power matrix. Smooth bounded $T$ therefore forces
$t\to0$ at $o$. Extending $t(o)=0$ gives a continuous nonnegative function
on compact $X$, nonzero because $T$ is normalized.
It has a positive maximum away from $o$.
The strong maximum principle \cite[pp.~131--132]{Gaud} on connected
$X\setminus\{o\}$ forces $t$ to be constant, contradicting its zero
limit at $o$.
\end{proof}

\section{The remaining probability questions}
Theorems~\ref{th:mass}, \ref{th:extension} and \ref{th:sublinear} give
restricted sufficient conditions. They leave the general attained problem
and its precise moment and entropy implications open.
For a smooth positive maximizing sequence set
\[
 d_j=D_p(T_j)\longrightarrow b^{-1/n},\qquad
 \nu_j=d_j^{-1}p k_{T_j}\,dV.
\]
Uniform first moments mean
$\int_X(-v)d\nu_j\leq C$ for every $j$ and every $h$-psh $v$ with
$\sup_Xv=0$. Uniform entropy means $\nu_j=a_jdV$ and
$\int_Xa_j\log a_j\,dV\leq C$ for every $j$.
Neither precise condition is proved here to imply \eqref{eq:G},
\eqref{eq:sublinear}, the smooth-root/rank hypotheses, or the desired
envelope for arbitrary $h$. Proposition~\ref{prop:excluded} fails the
attainment hypothesis and supplies no counterexample to either implication.

\begin{thebibliography}{9}
\bibitem{BEGZ} S. Boucksom, P. Eyssidieux, V. Guedj and A. Zeriahi,
Monge--Amp\`ere equations in big cohomology classes,
\emph{Acta Math.} \textbf{205} (2010), 199--262.
\bibitem{Dem} J.-P. Demailly, \emph{Complex Analytic and Differential Geometry},
author manuscript, 21 June 2012;
I.5.4--5.7, II.5.6, III.2.10--2.15, III.5.7--5.8,
III.6.8--6.9 and V.13.1--13.3;
\url{https://people.math.harvard.edu/~demarco/Math274/Demailly_ComplexAnalyticDiffGeom.pdf}.
\bibitem{Gaud} P. Gauduchon,
Les m\'etriques standard d'une surface complexe compacte
\`a premier nombre de Betti pair,
\emph{Ast\'erisque} \textbf{126} (1985), 129--135;
\url{https://www.numdam.org/article/AST_1985__126__129_0.pdf}.
\bibitem{GL} V. Guedj and C. H. Lu,
Quasi-plurisubharmonic envelopes 2: Bounds on Monge--Amp\`ere volumes,
\emph{Algebraic Geometry} \textbf{9} (2022), 688--713;
proof assessment uses arXiv:2106.04272, version 2 (2 March 2023).
\end{thebibliography}
\end{document}
